\chapter{Conclusion}
\label{chap:conclusion}

The conducted comparative analysis demonstrates that modern programming languages implement object-oriented principles through fundamentally different architectural approaches, each prioritizing a distinct balance between flexibility, formal safety, extensibility, runtime dynamism, and implementation simplicity. The calculated metrics revealed that no analyzed language provides a universally balanced object model across all considered dimensions. Languages oriented toward industrial stability and static correctness demonstrate stronger inheritance control and subtype formalization, while dynamically oriented systems provide higher flexibility of object representation and metaprogramming at the cost of weaker encapsulation and verification guarantees. At the same time, composition-oriented models reduce the structural complexity of inheritance hierarchies, but often sacrifice expressive polymorphism and formal abstraction mechanisms. The obtained results confirm that many classical problems of object-oriented programming are not caused by OOP itself, but rather by inconsistencies between inheritance semantics, subtype relations, mutability models, and runtime extensibility mechanisms.

The analysis also demonstrates several common architectural tendencies in the evolution of modern object systems. First, there is a gradual transition from inheritance-centered reuse toward composition-based models using interfaces, traits, mixins, and delegation. Second, contemporary languages increasingly attempt to separate subtype polymorphism from implementation inheritance, confirming the theoretical distinction established in type-system research. Third, there is a growing emphasis on deterministic hierarchy semantics, explicit override control, and restrictions intended to mitigate fragile base class effects. At the same time, the study revealed that mechanisms related to formal verification, immutability guarantees, concurrency safety, and contract-based specification remain insufficiently integrated into the majority of modern object-oriented languages. In most cases, such guarantees rely on external libraries, conventions, or runtime discipline rather than on the core semantics of the language itself.

Based on the identified limitations and recurring design trade-offs, this thesis proposes the conceptual foundation for a unified object model intended to combine strict semantic separation of inheritance and subtyping, composition-oriented reuse, deterministic polymorphic behavior, controlled metaprogramming, and stronger guarantees of state safety and formal verifiability. The proposed direction aims to preserve the expressive flexibility of object-oriented programming while reducing architectural ambiguity and improving correctness guarantees at the language level. Consequently, the results of this work may serve both as a systematic framework for further comparative studies of object-oriented languages and as a theoretical basis for the development of future programming languages with more balanced and formally robust object models.
